% FILE: 02_lineage_and_context.tex
% Project: experiments-visualizer-nextjs
% Role: web-based-pi-visualization-surface

\section{Lineage and Context}
\label{sec:experiments-visualizer-nextjs-lineage}

\subsection{2016 Language-Model Lesson}
\label{sec:experiments-visualizer-nextjs-lineage-2016}

The 2016 language-model experiment established that local text imitation does not conserve
consequence: a model can reproduce the surface form of a process description without preserving
the causal and temporal structure that makes that description meaningful as a claim about an
enterprise system. The lesson was that \emph{appearance} and \emph{evidence} are not the same
thing. A chart that looks like a conformance report is not a conformance report unless its
values are derived from a lawful replay of real event data against a verified process model.

The visualizer-nextjs project is a direct descendant of this lesson. The fact that its
\texttt{globals.css} defines a visually sophisticated process-intelligence dashboard --- Petri
net panels, alignment block sequences, EWMA alarm banners, cryptographic ledger explorers ---
while \texttt{page.tsx} remains the stock \texttt{create-next-app} welcome page is a precise
re-instantiation of the 2016 lesson at the UI layer: \emph{the surface vocabulary is present;
the causal connection to real evidence is absent.} The project is, at status PARTIAL, a
demonstration that a design system alone cannot constitute a visualization surface any more
than a text model alone can constitute a process-aware reasoning system.

\subsection{Enterprise Process Gap}
\label{sec:experiments-visualizer-nextjs-lineage-enterprise-gap}

Enterprise process intelligence systems fail for a predictable structural reason: the artefacts
they manufacture --- BPMN models, conformance reports, KPI dashboards --- are produced by
specialists in isolated toolchains and are never connected to the live event logs that would
allow them to be verified. Process mining tools exist in one silo; business intelligence
dashboards exist in another; neither is connected to the other by a lawful evidence chain.

The visualizer-nextjs project is designed to close this gap at the observation layer. It is
intended to be the single surface on which a researcher or auditor can see, simultaneously,
the Petri net model being replayed, the alignment moves that map the event log to the model,
the drift signal indicating whether the process is shifting, and the cryptographic receipt
chain that proves all upstream computation was honest. In the full architecture this surface
eliminates the need for specialist tooling by making process evidence browser-native and
auditable without command-line access. AUTHOR THESIS: this architectural position --- a
single web surface integrating all evidence classes --- is the correct response to the
enterprise silo failure mode.

\subsection{Relationship to the Chatman Equation}
\label{sec:experiments-visualizer-nextjs-lineage-chatman-eq}

The Chatman Equation posits that auditable process intelligence is a function of the
organization's observable event history: $A = \mu(O^{*})$, where $A$ is the admitted process
evidence, $\mu$ is the conformance-scoring operator, and $O^{*}$ is the complete object-centric
event log. The visualizer-nextjs project is the surface through which the output of $\mu$ is
rendered human-inspectable.

Concretely: the CSS custom properties in \texttt{globals.css} encode the value space of $A$.
The color tokens \texttt{--color-success}, \texttt{--color-warning}, \texttt{--color-danger},
and \texttt{--color-glow-green} / \texttt{--color-glow-red} correspond to the ordinal fitness
levels that $\mu$ can return. The \texttt{.ledger-block-card.healthy} and
\texttt{.ledger-block-card.tampered} states correspond to the binary integrity verdict of the
cryptographic receipt chain. Every CSS class in the design system is, in this framing, a
rendering rule for a term in the value algebra of $A = \mu(O^{*})$.

\subsection{Relationship to Receipts and Replay}
\label{sec:experiments-visualizer-nextjs-lineage-receipts}

The receipt and replay subsystem is the epistemic backbone of the research program. As
documented in the parent experiments corpus (\texttt{EVIDENCE\_CHAIN\_TRACE.md},
\texttt{checkpoint\_\_experiments\_complete.md}), the full evidence chain runs:
Raw~$\to$~Parsed~$\to$~Admitted~$\to$~Receipt. The visualizer-nextjs project is designed to
expose all four stages of this chain through a single browser session.

The \texttt{.receipt-json-body} and \texttt{.verification-console} CSS classes in
\texttt{globals.css} are the surface markers for the Receipt stage: they define how a receipt
JSON blob is displayed and how a verification verdict (pass or fail) is communicated to the
viewer. The \texttt{.ledger-chain-arrow.broken} class encodes the tamper-detection state ---
the visual signal that a hash-chain break has been detected, meaning one or more upstream
receipts cannot be verified. These classes constitute a specification for a receipt-aware
visualization layer that does not yet have a React implementation behind it.

The relationship is therefore one of specification-without-implementation: the receipt and
replay machinery exists upstream (verified by the EXPERIMENTS\_COMPLETE checkpoint dated
2026-05-31); the visualization surface has specified how to display its outputs but has not
yet wired up the connection.
